% Generated draft: requires mathematical and source review.
\documentclass[11pt]{article}
\usepackage[margin=25mm]{geometry}
\usepackage{fontspec}
\setmainfont{Latin Modern Roman}
\usepackage{amsmath,amssymb,mathtools,unicode-math}
\setmathfont{Latin Modern Math}
\usepackage{enumitem,needspace,xurl,hyperref}
\setlength{\parindent}{0pt}
\setlength{\parskip}{5pt}
\setlength{\emergencystretch}{3em}
\newcommand{\N}{\mathbb{N}}
\newcommand{\Z}{\mathbb{Z}}
\newcommand{\Q}{\mathbb{Q}}
\newcommand{\R}{\mathbb{R}}
\newcommand{\C}{\mathbb{C}}
\newcommand{\F}{\mathbb{F}}
\newcommand{\A}{\mathbb{A}}
\newcommand{\PP}{\mathbb{P}}
\newcommand{\E}{\mathbb{E}}
\newcommand{\eps}{\varepsilon}
\newcommand{\dd}{\,\mathrm{d}}
\begin{document}
\section*{UnsolvedMath ID: 20002334}
\textbf{Source problem code:} catalogue entry
\subsection{Definitions and mathematical statement}
The source-owned item is Problem 20: ``Study non commutative probability spaces,'' with the note that the workshop's Non Commutative Probability Spaces working group studied this subject. This is a research programme, not a single quantified conjecture. A noncommutative probability space may mean a unital \(*\)-algebra with a positive normalized linear functional. The supplied partial result makes the narrower tracial von Neumann algebra assumptions explicit below; it does not impose them retroactively on every possible interpretation of Problem 20.

Let \((M,\tau)\) and \((N,\sigma)\) be von Neumann algebras with faithful normal tracial states, normalized by \(\tau(1)=\sigma(1)=1\). Write \(\|x\|_{2,\tau}=\tau(x^*x)^{1/2}\), and similarly for \(\sigma\). Let \(a=(a_1,\ldots,a_r)\in M^r\) and \(b=(b_1,\ldots,b_r)\in N^r\) be tuples of operator-norm contractions. A joint \(*\)-moment is \(\tau(w(a,a^*))\), for a word \(w\) in the formal variables \(X_1,\ldots,X_r,X_1^*,\ldots,X_r^*\), including the empty word.

The aid asserts the equivalence of the following three conditions:
\[
\begin{aligned}
 &\tau(w(a,a^*))=\sigma(w(b,b^*))\quad\text{for every word }w;\\
 &\exists\text{ a normal trace-preserving }*\text{-isomorphism }
 \theta:W^*(a_1,\ldots,a_r)\to W^*(b_1,\ldots,b_r),
 \quad\theta(a_i)=b_i;\\
 &\operatorname{qftp}^{M}(a)=\operatorname{qftp}^{N}(b)
 \quad\text{in the usual continuous tracial von Neumann algebra language}.
\end{aligned}
\]
Here \(W^*(a)\) is the generated von Neumann algebra, and quantifier-free type means the values of all quantifier-free formulas on the tuple, using the \(2\)-norm metric and the real and imaginary trace predicates in the standard operator-norm-ball sorts. The isomorphism is of the generated algebras, not necessarily the whole ambient \(M\) and \(N\). No equality of full types with quantifiers is asserted.

The quantitative statement is as follows. Fix an integer \(d\ge0\) and \(\varepsilon\ge0\), and assume
\[
 |\tau(w(a,a^*))-\sigma(w(b,b^*))|\le\varepsilon
 \qquad(|w|\le2d).
\]
For every noncommutative \(*\)-polynomial \(p=\sum_{|w|\le d}c_w w\), put \(\|c(p)\|_1=\sum_w|c_w|\). Then
\[
 \left|\|p(a)\|_{2,\tau}^{,2}-\|p(b)\|_{2,\sigma}^{,2}\right|
 \le\|c(p)\|_1^2\varepsilon,
 \qquad
 \left|\|p(a)\|_{2,\tau}-\|p(b)\|_{2,\sigma}\right|
 \le\|c(p)\|_1\sqrt{\varepsilon}.
\]
Trace atoms also satisfy
\[
 |\tau(p(a))-\sigma(p(b))|\le\|c(p)\|_1\varepsilon
\]
when their degrees lie within the controlled moment range. Consequently, for a prescribed quantifier-free formula, choose \(d\) large enough for its finitely many polynomial atoms and propagate the atomic error bounds through the specified continuity moduli of its connectives. If a connective \(F\) has modulus \(\omega_F\) for the maximum metric on its inputs and their errors are at most \(\eta\), its output error is at most \(\omega_F(\eta)\). An effective numerical certificate requires effective or explicitly supplied moduli; arbitrary unspecified continuous functions are not automatically computable. These statements describe one partial approach to Problem 20.

The registry concatenates later independent items, which are excluded from this target: Problem 21 asks for NIP continuous theories and whether randomization preserves NIP; Problem 22 concerns stable and \(\omega\)-stable groups; Problem 23 concerns definable sets; Problem 24 discusses two-sorted weak structures \((M,[0,1],\ldots)\), symbols for continuous functions, standard-part metric structures, the use of \(\sup\) and \(\inf\), and Keisler--Shelah and omitting-types theorems; Problem 25 asks about integrals versus types. The respective Stable Groups, Definable Sets, and Weak First Order and Probability Structures working-group background notes are contextual rather than hypotheses of the moment statement.
\subsection{Short English statement}
Problem 20 is a broad programme on noncommutative probability. In the supplied tracial von Neumann algebra setting, joint moments determine the generated algebra and quantifier-free type, and finite moment errors control polynomial \(2\)-norms and prescribed quantifier-free formulas.
\subsection{Sources}
\begin{itemize}
\item Local registry, permanent UnsolvedMath ID \(20002334\), exact public code \texttt{catalogue entry}; canonical location \texttt{aim-logic-notes.json}, \texttt{notes[109]}.
\item AIM, \emph{Questions on Model theory for metric structures}, Problem 20, page 4, and its Non Commutative Probability Spaces working-group background; supplied reference \url{https://aimath.org/WWN/continuouslogic/continuouslogic.pdf}. The source-verification field reports that the official document separates the later numbered items accidentally concatenated in the registry. No such inspection was performed here.
\item Supplied research attempt 1: \texttt{partial\_result}, run and registry status \texttt{partially\_solved}, rigor label \texttt{proved}, low novelty confidence; classification \texttt{PARTIAL-PROGRESS}, created and updated 14 August 2026. No proposer, year, original literature entry, or primary reference for the moment and quantitative statements is supplied.
\item The research aid explicitly states that it is machine generated, not peer reviewed, and requires independent verification before citation.
\end{itemize}
\end{document}
